\chapter{Objetivos}

\section{Objetivo geral}

Desenvolver e analisar uma sequência didática apoiada por uma plataforma interativa para o ensino introdutório de \textit{machine learning} como conteúdo de Computação no Ensino Médio, examinando mudanças e permanências nas explicações dos estudantes sobre como os dados e a finalidade da tarefa se relacionam com regras, parâmetros e resultados dos modelos.

\section{Objetivos específicos}

Os objetivos específicos compreendem:

\begin{itemize}
    \item analisar referências sobre o ensino de inteligência artificial e \textit{machine learning} voltadas à Educação Básica, relacionando o uso de representações visuais interativas à mediação docente e à avaliação formativa da aprendizagem;
    \item relacionar os objetivos de aprendizagem da sequência didática aos eixos e às habilidades pertinentes da Computação na Educação Básica, sem supor que as diretrizes curriculares vigentes prescrevam algoritmos específicos para a escola;
    \item organizar uma sequência didática sobre regressão linear, máquina de vetores de suporte (\textit{SVM}), \textit{k}-vizinhos mais próximos (\textit{k}-NN), árvore de decisão e \textit{k}-means, partindo de situações-problema que explicitem a origem dos dados para formular previsões, testar alterações de forma controlada e explicar os resultados observados;
    \item ajustar às condições reais do campo de aplicação tanto os controles visuais e as explicações matemáticas da plataforma quanto a construção da tabela inicial pela própria turma e o uso dos cenários sintéticos nos módulos;
    \item aplicar a sequência em uma turma na qual o pesquisador atua como professor;
    \item identificar o que muda ou permanece nas dificuldades dos estudantes para explicar como os dados e a finalidade da tarefa se relacionam com a alteração de regras ou parâmetros e com os resultados produzidos pelos modelos estudados, confrontando o diagnóstico inicial, os três registros intermediários e o instrumento final;
    \item examinar, diante de uma situação-problema inédita em relação às atividades trabalhadas nos encontros, como os estudantes mobilizam os conceitos estudados ao longo da sequência didática para justificar escolhas e interpretar os resultados;
    \item sistematizar como produto educacional a plataforma articulada à sequência didática e aos respectivos instrumentos.
\end{itemize}
